\chapter{Điểm Số Không Phải Tất Cả}
\markboth{Điểm Số Không Phải Tất Cả}{Scores Are Not Everything}

\section{Điểm Cao Không Tự Động Biến Em Thành Người Giỏi}
\begin{truyen}{Điểm Cao Không Tự Động Biến Em Thành Người Giỏi}{High Scores Do Not Automatically Make You Capable}
\chuhoa{M}{inh luôn đứng nhất lớp.}
\textit{(Minh was always first in class.)}

Bảng điểm của em gần như không có số nào dưới chín.
\textit{(His report card almost never showed a score below nine.)}

Nhưng mỗi khi thầy hỏi thêm một câu ngoài sách giáo khoa, Minh đứng im.
\textit{(But whenever the teacher asked one extra question beyond the textbook, Minh froze.)}

Em biết đáp án mẫu, nhưng không biết nghĩ tiếp.
\textit{(He knew the model answer, but did not know how to think beyond it.)}

Đến khi thi học sinh giỏi, đề không hỏi học thuộc.
\textit{(When the gifted-student exam came, the questions did not reward memorization.)}

Đề bắt phải hiểu.
\textit{(They demanded understanding.)}

Minh làm bài như người đi trên chiếc cầu đã mất bản đồ.
\textit{(Minh wrote like someone walking on a bridge after losing the map.)}

Điểm số từng giúp em được khen.
\textit{(Scores once brought him praise.)}

Nhưng chính điểm số cũng che giấu lỗ hổng lớn nhất: em học để đúng, không học để hiểu.
\textit{(But they also hid the largest gap: he studied to be correct, not to understand.)}

Thầy giáo nói thẳng với cả lớp: "Điểm cao là tín hiệu tốt, nhưng không phải bằng chứng cuối cùng của năng lực.
\textit{(The teacher told the class plainly: "High scores are a good signal, but not the final proof of competence.)}

Năng lực lộ ra khi em gặp câu chưa từng thấy."
\textit{(Competence shows when you face a question you have never seen before.")}

\end{truyen}

\begin{baihoc}
Đừng dùng điểm số để ru ngủ mình. Hãy tự hỏi: ngoài bài kiểm tra, mình giải thích lại được không, ứng dụng được không, và tự học tiếp được không?
\textit{(Do not let scores lull you to sleep. Ask instead: beyond the test, can I explain it, apply it, and continue learning on my own?)}
\end{baihoc}

\begin{ghinhoanh}
\textbf{Short English to remember:}

Do not let scores lull you to sleep.

Ask instead: beyond the test, can I explain it, apply it, and continue learning on my own?

\end{ghinhoanh}

\begin{tuhocnhanh}
\textbf{Cau hoi tu hoc / Self-study questions}

1. Van de chinh la gi? \textit{What is the main problem?}

2. Cach xu ly tot hon la gi? \textit{What is a better action?}

3. Mot cau tieng Anh em muon nho la cau nao? \textit{Which English sentence do you want to remember?}
\end{tuhocnhanh}
\ngancach

\section{Điểm Thấp Không Có Nghĩa Em Vô Dụng}
\begin{truyen}{Điểm Thấp Không Có Nghĩa Em Vô Dụng}{Low Scores Do Not Mean You Are Worthless}
\chuhoa{L}{an bị điểm bốn môn Toán.}
\textit{(Lan got a four in mathematics.)}

Em giấu bài kiểm tra xuống đáy cặp như giấu một vết thương.
\textit{(She hid the test at the bottom of her bag as if hiding a wound.)}

Tối đó em nói: ``Con ngu thật.'' Mẹ em không mắng, chỉ hỏi: ``Bài nào con sai vì không biết, bài nào con sai vì cẩu thả?''
\textit{(That evening she said: ``I am stupid.'' Her mother did not scold her, but asked: ``Which mistakes came from not knowing, and which came from carelessness?'')}

Hai mẹ con ngồi tách từng lỗi ra.
\textit{(Together they separated each mistake.)}

Hóa ra ba câu sai vì quên đổi đơn vị, hai câu sai vì đọc thiếu dữ kiện, chỉ một câu là thực sự không hiểu.
\textit{(It turned out three errors came from forgetting units, two from misreading the problem, and only one from truly not understanding.)}

Điểm bốn vẫn là điểm bốn.
\textit{(A four was still a four.)}

Nhưng từ chỗ thấy mình kém cỏi toàn bộ, Lan bắt đầu nhìn ra một sự thật công bằng hơn: mình sai ở đâu thì sửa ở đó.
\textit{(But instead of seeing herself as entirely inadequate, Lan began to see a fairer truth: fix the specific place where you are wrong.)}

Một tháng sau, Lan được bảy điểm.
\textit{(A month later, Lan scored seven.)}

Không phải phép màu.
\textit{(It was not magic.)}

Chỉ là em thôi chửi bản thân và bắt đầu sửa lỗi thật.
\textit{(She had simply stopped insulting herself and started correcting real errors.)}

\end{truyen}

\begin{baihoc}
Điểm thấp là dữ liệu, không phải bản án. Người tiến bộ không phải người chưa từng sai, mà là người chịu nhìn thẳng lỗi sai mà không tự hủy mình.
\textit{(A low score is data, not a verdict. The person who improves is not the one who never errs, but the one who can face mistakes without destroying self-respect.)}
\end{baihoc}

\begin{ghinhoanh}
\textbf{Short English to remember:}

A low score is data, not a verdict.

The person who improves is not the one who never errs, but the one who can face mistakes without destroying self-respect.

\end{ghinhoanh}

\begin{tuhocnhanh}
\textbf{Cau hoi tu hoc / Self-study questions}

1. Van de chinh la gi? \textit{What is the main problem?}

2. Cach xu ly tot hon la gi? \textit{What is a better action?}

3. Mot cau tieng Anh em muon nho la cau nao? \textit{Which English sentence do you want to remember?}
\end{tuhocnhanh}
\ngancach

\section{Có Những Đứa Trẻ Giỏi Vì Sợ}
\begin{truyen}{Có Những Đứa Trẻ Giỏi Vì Sợ}{Some Children Perform Well Because They Are Afraid}
\chuhoa{V}{y học ngày học đêm.}
\textit{(Vy studied day and night.)}

Ai nhìn vào cũng nghĩ em có ý chí thép.
\textit{(Anyone watching thought she had iron discipline.)}

Chỉ có em biết thứ kéo mình chạy không phải đam mê, mà là sợ hãi.
\textit{(Only she knew that what drove her was not passion, but fear.)}

Em sợ thua bạn, sợ bố thất vọng, sợ mình không còn giá trị nếu không đứng đầu.
\textit{(She feared losing to classmates, feared disappointing her father, feared becoming worthless if she was not first.)}

Thành tích của em tăng, nhưng giấc ngủ của em giảm.
\textit{(Her achievements rose, but her sleep shrank.)}

Nét mặt em căng như dây đàn.
\textit{(Her face stayed tight like a string.)}

Một đứa trẻ được gọi là giỏi, nhưng sống như đang chạy trốn.
\textit{(She was called gifted, yet lived as if fleeing something.)}

Cô chủ nhiệm nói riêng với em: "Nếu thứ giữ em học là nỗi sợ, rồi sẽ đến lúc em ghét luôn cả việc học.
\textit{(Her homeroom teacher told her privately: "If fear is what keeps you studying, one day you will hate learning itself.)}

Không ai chạy mãi bằng hoảng loạn được."
\textit{(No one can run forever on panic.")}

Vy bắt đầu nghỉ đúng giờ hơn, chừa một buổi tối không học thêm, và tập trả lời câu hỏi khó nhất: nếu không đứng nhất, mình còn là ai?
\textit{(Vy began ending study at a fixed hour, keeping one evening free from extra classes, and practicing the hardest question: if I am not first, who am I still?)}

\end{truyen}

\begin{baihoc}
Kết quả tốt đạt được bằng hoảng loạn không bền. Hãy xây động lực từ ý nghĩa, sự tò mò và lòng tự trọng, không phải từ nỗi sợ bị bỏ rơi hay bị so sánh.
\textit{(Results achieved through panic do not last. Build motivation from meaning, curiosity, and self-respect, not from fear of abandonment or comparison.)}
\end{baihoc}

\begin{ghinhoanh}
\textbf{Short English to remember:}

Results achieved through panic do not last.

Build motivation from meaning, curiosity, and self-respect, not from fear of abandonment or comparison.

\end{ghinhoanh}

\begin{tuhocnhanh}
\textbf{Cau hoi tu hoc / Self-study questions}

1. Van de chinh la gi? \textit{What is the main problem?}

2. Cach xu ly tot hon la gi? \textit{What is a better action?}

3. Mot cau tieng Anh em muon nho la cau nao? \textit{Which English sentence do you want to remember?}
\end{tuhocnhanh}
\ngancach

\section{Câu Hỏi Quan Trọng Hơn Điểm Số}
\begin{truyen}{Câu Hỏi Quan Trọng Hơn Điểm Số}{The Questions Matter More Than the Score}
\chuhoa{M}{ột thầy giáo mới vào lớp và viết lên bảng: ``Em muốn trở thành kiểu người như thế nào?''}
\textit{(A new teacher entered the class and wrote on the board: ``What kind of person do you want to become?'')}

Cả lớp ngơ ngác.
\textit{(The class stared blankly.)}

Họ quen với những câu hỏi như công thức nào, định lý nào, tác giả nào.
\textit{(They were used to questions like which formula, which theorem, which author.)}

Không ai quen bị hỏi về chính mình.
\textit{(No one was used to being asked about themselves.)}

Thầy nói: "Nếu em chỉ hỏi làm sao được chín, mười, em sẽ trở thành cỗ máy săn điểm.
\textit{(The teacher said: "If you only ask how to get nines and tens, you will become a score-hunting machine.)}

Nếu em hỏi làm sao hiểu sâu hơn, sống tử tế hơn, làm việc chắc hơn, em mới đang học thật."
\textit{(If you ask how to understand more deeply, live more decently, and work more solidly, then you are truly learning.")}

Buổi học hôm đó không có nhiều bài tập.
\textit{(There were not many exercises that day.)}

Nhưng nhiều học sinh về nhà với một cảm giác lạ: lần đầu tiên việc học được nối với cuộc đời thật.
\textit{(But many students went home with a strange feeling: for the first time, study had been connected to real life.)}

\end{truyen}

\begin{baihoc}
Điểm số quan trọng, nhưng câu hỏi em mang theo suốt đời còn quan trọng hơn. Học để thành người, rồi mới học để thắng bài kiểm tra.
\textit{(Scores matter, but the questions you carry for life matter more. Study first to become a person, then to beat a test.)}
\end{baihoc}

\begin{ghinhoanh}
\textbf{Short English to remember:}

Scores matter, but the questions you carry for life matter more.

Study first to become a person, then to beat a test.

\end{ghinhoanh}

\begin{tuhocnhanh}
\textbf{Cau hoi tu hoc / Self-study questions}

1. Van de chinh la gi? \textit{What is the main problem?}

2. Cach xu ly tot hon la gi? \textit{What is a better action?}

3. Mot cau tieng Anh em muon nho la cau nao? \textit{Which English sentence do you want to remember?}
\end{tuhocnhanh}
\ngancach

\section{Sĩ Tử Đèn Dầu Và Cái Đầu Rỗng}
\begin{truyen}{Sĩ Tử Đèn Dầu Và Cái Đầu Rỗng}{The Lamp-Lit Candidate With an Empty Head}
\chuhoa{N}{gày xưa có một sĩ tử nổi tiếng chăm chỉ.}
\textit{(Long ago there was a candidate famous for diligence.)}

Đêm nào trong làng người ta cũng thấy ánh đèn dầu nhà cậu sáng tới khuya.
\textit{(Every night the villagers saw his oil lamp burning late.)}

Ai cũng tin cậu sẽ đỗ cao.
\textit{(Everyone believed he would rank highly.)}

Nhưng vào trường thi, đề hỏi một ý lạ hơn sách cũ một chút là cậu lúng túng ngay.
\textit{(But in the examination hall, the moment the question moved slightly beyond the standard text, he was lost.)}

Thầy đồ già nói: "Con đốt nhiều dầu thật, nhưng chưa chắc đã soi sáng đầu mình.
\textit{(The old teacher said, "You have burned much oil, but not necessarily lit up your mind.)}

Học mà chỉ nhét chữ thì đầu nặng chứ não chưa mở."
\textit{(If study is only stuffing words in, the head becomes heavy while the mind stays shut.")}

Từ đó cậu đổi cách học: bớt chép, tăng hỏi, tăng giải thích lại bằng lời của mình.
\textit{(From then on he changed his method: less copying, more questioning, more explaining in his own words.)}

Năm sau cậu mới thi đỗ thật.
\textit{(Only the next year did he truly pass.)}

\end{truyen}

\begin{baihoc}
Thời nào cũng vậy, học nhiều giờ không đồng nghĩa học sâu. Đừng nhầm số giờ với chất lượng tư duy.
\textit{(In every era, studying for many hours does not equal studying deeply. Do not confuse the number of hours with the quality of thought.)}
\end{baihoc}

\begin{ghinhoanh}
\textbf{Short English to remember:}

In every era, studying for many hours does not equal studying deeply.

Do not confuse the number of hours with the quality of thought.

\end{ghinhoanh}

\begin{tuhocnhanh}
\textbf{Cau hoi tu hoc / Self-study questions}

1. Van de chinh la gi? \textit{What is the main problem?}

2. Cach xu ly tot hon la gi? \textit{What is a better action?}

3. Mot cau tieng Anh em muon nho la cau nao? \textit{Which English sentence do you want to remember?}
\end{tuhocnhanh}
\ngancach

\section{Học Tủ Là Một Canh Bạc}
\begin{truyen}{Học Tủ Là Một Canh Bạc}{Selective Studying Is a Gamble}
\chuhoa{T}{ùng nghe anh khóa trên mách rằng đề năm nào cũng quanh quẩn vài dạng.}
\textit{(Tung heard older students say that the exam always repeated a few patterns.)}

Em bỏ hẳn mấy phần khó và chỉ học đúng những gì mình tin sẽ ra.
\textit{(He dropped the hard parts and studied only what he believed would appear.)}

Tuần trước khi thi, em còn rất tự tin vì làm đề đoán khá trúng.
\textit{(A week before the exam, he was confident because the predicted papers fit him well.)}

Đến hôm thi thật, đề đổi hướng.
\textit{(On the real exam day, the paper shifted direction.)}

Tùng ngồi nhìn giấy như nhìn một người lạ.
\textit{(Tung stared at it as if at a stranger.)}

Em không trượt vì kém thông minh.
\textit{(He did not fail because he lacked intelligence.)}

Em trượt vì đã biến việc học thành trò đặt cược.
\textit{(He failed because he had turned learning into a betting game.)}

\end{truyen}

\begin{baihoc}
Học tủ có thể thắng một lần, nhưng rất khó nuôi năng lực thật. Điều em cần là nền rộng đủ để không sập khi đề đổi.
\textit{(Selective studying may win once, but it rarely builds real ability. What you need is a foundation wide enough not to collapse when the exam changes.)}
\end{baihoc}

\begin{ghinhoanh}
\textbf{Short English to remember:}

Selective studying may win once, but it rarely builds real ability.

What you need is a foundation wide enough not to collapse when the exam changes.

\end{ghinhoanh}

\begin{tuhocnhanh}
\textbf{Cau hoi tu hoc / Self-study questions}

1. Van de chinh la gi? \textit{What is the main problem?}

2. Cach xu ly tot hon la gi? \textit{What is a better action?}

3. Mot cau tieng Anh em muon nho la cau nao? \textit{Which English sentence do you want to remember?}
\end{tuhocnhanh}
\ngancach

\section{Cô Bé Bán Hàng Tính Nhanh Hơn Học Sinh Giỏi}
\begin{truyen}{Cô Bé Bán Hàng Tính Nhanh Hơn Học Sinh Giỏi}{The Shop Girl Who Calculated Faster Than the Top Student}
\chuhoa{Ở}{ chợ, một cô bé phụ mẹ bán rau tính tiền rất nhanh.}
\textit{(At the market, a girl helping her mother sell vegetables calculated very quickly.)}

Khách vừa gọi ba món là em đã nhẩm ra số tiền và tiền thối gần như ngay lập tức.
\textit{(The moment customers named three items, she almost instantly knew the total and the change.)}

Một học sinh giỏi Toán đi ngang qua ngạc nhiên.
\textit{(A top mathematics student passing by was surprised.)}

Em ấy giải phương trình tốt hơn nhiều, nhưng gặp tình huống thật lại chậm hơn cô bé bán hàng.
\textit{(He could solve equations far better, but in real situations he was slower than the girl at the stall.)}

Người mẹ cười: "Mỗi người giỏi một kiểu.
\textit{(The mother laughed: "People are good in different ways.)}

Học trong trường là một chuyện, dùng được ngoài đời là chuyện khác nữa."
\textit{(Learning in school is one thing; using it in life is another.")}

Cậu học sinh nghe xong thấy xấu hổ theo nghĩa tốt.
\textit{(The boy felt a useful kind of embarrassment.)}

Lần đầu tiên em hiểu rằng kiến thức không đi ra được đời thì vẫn còn thiếu một đoạn đường.
\textit{(For the first time he understood that knowledge which cannot step into life is still incomplete.)}

\end{truyen}

\begin{baihoc}
Trường học không phải toàn bộ năng lực con người. Điểm cao đáng quý, nhưng khả năng làm được việc thật cũng đáng quý không kém.
\textit{(School is not the whole of human ability. High scores are valuable, but the ability to do real things is equally valuable.)}
\end{baihoc}

\begin{ghinhoanh}
\textbf{Short English to remember:}

School is not the whole of human ability.

High scores are valuable, but the ability to do real things is equally valuable.

\end{ghinhoanh}

\begin{tuhocnhanh}
\textbf{Cau hoi tu hoc / Self-study questions}

1. Van de chinh la gi? \textit{What is the main problem?}

2. Cach xu ly tot hon la gi? \textit{What is a better action?}

3. Mot cau tieng Anh em muon nho la cau nao? \textit{Which English sentence do you want to remember?}
\end{tuhocnhanh}
\ngancach

\section{Bài Kiểm Tra Mở Sách}
\begin{truyen}{Bài Kiểm Tra Mở Sách}{The Open-Book Test}
\chuhoa{M}{ột cô giáo bất ngờ cho lớp làm bài kiểm tra mở sách.}
\textit{(A teacher surprised the class with an open-book test.)}

Cả lớp mừng vì tưởng sẽ dễ.
\textit{(Everyone was happy because they thought it would be easy.)}

Nhưng khi nhận đề, nhiều em hoảng hơn bình thường.
\textit{(But when they received the paper, many were more panicked than usual.)}

Sách ở đó, nhưng câu hỏi buộc phải biết nối ý, so sánh, và tự kết luận.
\textit{(The book was there, but the questions required connecting ideas, comparing, and drawing conclusions.)}

Cô giáo nói: "Thi đóng sách kiểm tra trí nhớ.
\textit{(The teacher said, "A closed-book test checks memory.)}

Thi mở sách kiểm tra đầu óc các em dùng trí nhớ vào việc gì."
\textit{(An open-book test checks what your mind does with memory.")}

Buổi đó nhiều học sinh hiểu ra thứ mình còn yếu nhất không phải nhớ ít, mà là nghĩ ít.
\textit{(That day many students realized their greatest weakness was not remembering too little, but thinking too little.)}

\end{truyen}

\begin{baihoc}
Đừng chỉ luyện nhớ. Hãy luyện kết nối, lập luận và giải thích. Đó mới là phần khó bị thay thế nhất.
\textit{(Do not train memory alone. Train connection, reasoning, and explanation. Those are the parts hardest to replace.)}
\end{baihoc}

\begin{ghinhoanh}
\textbf{Short English to remember:}

Do not train memory alone.

Train connection, reasoning, and explanation.

Those are the parts hardest to replace.

\end{ghinhoanh}

\begin{tuhocnhanh}
\textbf{Cau hoi tu hoc / Self-study questions}

1. Van de chinh la gi? \textit{What is the main problem?}

2. Cach xu ly tot hon la gi? \textit{What is a better action?}

3. Mot cau tieng Anh em muon nho la cau nao? \textit{Which English sentence do you want to remember?}
\end{tuhocnhanh}
\ngancach

\section{Khi Điểm Mười Là Mục Tiêu Duy Nhất}
\begin{truyen}{Khi Điểm Mười Là Mục Tiêu Duy Nhất}{When a Perfect Score Becomes the Only Goal}
\chuhoa{C}{ó học sinh làm bài được chín rưỡi nhưng chỉ chăm chăm hỏi: ``Sao em mất nửa điểm?''}
\textit{(Some students score nine and a half yet ask only, ``Why did I lose half a point?'')}

Câu hỏi đó không sai, nhưng nếu nó là câu duy nhất, việc học sẽ bị thu nhỏ thành săn lỗi để đủ điểm tối đa.
\textit{(That question is not wrong, but if it is the only question, learning shrinks into a hunt for lost marks.)}

Một thầy giáo trả bài rồi hỏi ngược: ``Ngoài nửa điểm đó, em đã hiểu sâu thêm điều gì?'' Học sinh đứng im vì chưa từng nghĩ tới.
\textit{(A teacher handed back the paper and asked in return, ``Besides that half point, what did you understand more deeply?'' The student stood silent, never having considered it.)}

Nhiều em không thiếu năng lực.
\textit{(Many students do not lack ability.)}

Các em chỉ bị hệ thống phần thưởng làm cho quên mất mục đích thật của việc học.
\textit{(They have simply been shaped by reward systems that make them forget the real purpose of study.)}

\end{truyen}

\begin{baihoc}
Điểm tối đa không xấu. Nhưng nếu nó là mục tiêu duy nhất, em rất dễ trở thành người học hẹp.
\textit{(A perfect score is not bad. But if it becomes the only goal, you easily become a narrow learner.)}
\end{baihoc}

\begin{ghinhoanh}
\textbf{Short English to remember:}

A perfect score is not bad.

But if it becomes the only goal, you easily become a narrow learner.

\end{ghinhoanh}

\begin{tuhocnhanh}
\textbf{Cau hoi tu hoc / Self-study questions}

1. Van de chinh la gi? \textit{What is the main problem?}

2. Cach xu ly tot hon la gi? \textit{What is a better action?}

3. Mot cau tieng Anh em muon nho la cau nao? \textit{Which English sentence do you want to remember?}
\end{tuhocnhanh}
\ngancach

\section{Học Để Làm Việc Hay Học Để Diễn?}
\begin{truyen}{Học Để Làm Việc Hay Học Để Diễn?}{Are You Studying to Work or to Perform?}
\chuhoa{M}{ột nhóm học sinh chuẩn bị thuyết trình.}
\textit{(A group of students prepared a presentation.)}

Có bạn chăm sửa slide cho thật đẹp, có bạn lo tập nói, có bạn đọc thêm tài liệu để hiểu vấn đề.
\textit{(One polished the slides, one practiced delivery, one read further to truly understand the topic.)}

Hôm trình bày, hình thức của nhóm rất ổn.
\textit{(On presentation day, the group looked polished.)}

Nhưng chỉ cần giáo viên hỏi sâu hơn một chút, cả nhóm lộ ra chỗ rỗng ngay.
\textit{(But as soon as the teacher asked a slightly deeper question, the emptiness showed.)}

Có nhiều thứ trong trường học cho phép em diễn được khá lâu.
\textit{(Many things in school allow you to perform convincingly for quite a while.)}

Nhưng đời thật sẽ sớm hỏi tới phần ruột.
\textit{(Real life, however, quickly asks about substance.)}

Sau buổi đó, cả nhóm mới hiểu: học không chỉ để qua mặt một bài chấm.
\textit{(After that session, the group finally understood: studying is not only to get past a grade.)}

Học là để đứng được khi bị hỏi thật.
\textit{(It is to remain standing when real questions come.)}

\end{truyen}

\begin{baihoc}
Nếu em chỉ trau vỏ ngoài, em sẽ sống trong nỗi sợ bị hỏi thật. Năng lực thật là thứ làm em bớt sợ.
\textit{(If you polish only the surface, you will live in fear of real questions. Real competence is what makes that fear smaller.)}
\end{baihoc}

\begin{ghinhoanh}
\textbf{Short English to remember:}

If you polish only the surface, you will live in fear of real questions.

Real competence is what makes that fear smaller.

\end{ghinhoanh}

\begin{tuhocnhanh}
\textbf{Cau hoi tu hoc / Self-study questions}

1. Van de chinh la gi? \textit{What is the main problem?}

2. Cach xu ly tot hon la gi? \textit{What is a better action?}

3. Mot cau tieng Anh em muon nho la cau nao? \textit{Which English sentence do you want to remember?}
\end{tuhocnhanh}
